\documentclass[lettersize,journal]{IEEEtran}
\usepackage{amsmath,amsfonts}
\usepackage{algorithmic}
\usepackage{algorithm}
\usepackage{array}
\usepackage[caption=false,font=normalsize,labelfont=sf,textfont=sf]{subfig}
\usepackage{textcomp}
\usepackage{stfloats}
\usepackage{url}
\usepackage{verbatim}
\usepackage{graphicx}
\usepackage{multirow}
\usepackage{booktabs}
\usepackage{xcolor}
\usepackage{cite}
\hyphenation{op-tical net-works semi-conduc-tor IEEE-Xplore}
% updated with editorial comments 8/9/2021

\begin{document}

\title{TinySplat: Feedforward Approach for Generating Compact 3D Scene Representation}

\author{
	Zetian~Song,
        Jiaye~Fu,
	Jiaqi~Zhang,
        Xiaohan~Lu,\\
        Chuanmin~Jia, ~\IEEEmembership{Member,~IEEE,}
	Siwei~Ma, ~\IEEEmembership{Fellow,~IEEE,}
	Wen~Gao, ~\IEEEmembership{Fellow,~IEEE,}\\	
		\thanks{
			%This work was supported by the National Key Research and Development Project (2019YFF0302703), National Natural Science Foundation of China (61931014, 62031013), and in part by the High Performance Computing Platform of Peking University, which are gratefully acknowledged. (Corresponding author: Jiaqi Zhang.)
		}
		% \thanks{Zetian~Song is with the Shenzhen International Graduate School, Tsinghua University, Shenzhen 518055, China and also with the Konka Grp Co Ltd, Shenzhen 518057, China~(e-mail: guobin@konka.com)}
		% \thanks{Jiaye~Fu is with the Konka Grp Co Ltd, Shenzhen 518057, China~(e-mail: fuyuhong@konka.com)}
		% \thanks{Xiaohan~Lu is with the Shenzhen International Graduate School, Tsinghua University, Shenzhen 518055, China~(e-mail: mahui@tsinghua.edu.cn)}
		% \thanks{Jiaqi~Zhang is with the BEIJING BOYA RUISHI TECHNOLOGY CO., LTD, Beijing 100085, China~(email: jcwei@realscene.cn)}
        \thanks{Zetian~Song, Jiaqi~Zhang, Xiaohan~Lu, Siwei~Ma and Wen Gao are with the State Key Laboratory of Multimedia Information Processing, School of Computer Science, Peking University, Beijing, 100871, China~(e-mail: songzt@pku.edu.cn, jqzhang@pku.edu.cn, luxiaohan@stu.pku.edu.cn, swma@pku.edu.cn, wgao@pku.edu.cn). }
        \thanks{Jiaye~Fu is with the School of Electronic and Computer Engineering, Peking University, Shenzhen, 518055, China, and also with the State Key Laboratory of Multimedia Information Processing, School of Computer Science, Peking University, Beijing, 100871, China (email: jyfu@stu.pku.edu.cn).}
        \thanks{Chuanmin~Jia is with the Wangxuan Institute of Computer Technology, State Key Laboratory of Multimedia Information Processing, Peking University, Beijing 100871, China~(email: cmjia@pku.edu.cn).}
	}

% The paper headers
\markboth{IEEE Transactions on Circuits and Systems for Video Technology}%
{Shell \MakeLowercase{\textit{et al.}}: A Sample Article Using IEEEtran.cls for IEEE Journals}

% \IEEEpubid{0000--0000/00\$00.00~\copyright~2021 IEEE}
% Remember, if you use this you must call \IEEEpubidadjcol in the second
% column for its text to clear the IEEEpubid mark.

\maketitle

\begin{abstract}
The recent development of feedforward 3D Gaussian Splatting~(3DGS) presents a new paradigm to reconstruct 3D scenes. Using neural networks trained on large-scale multi-view datasets, it can directly infer 3DGS representations from sparse input views. Although the feedforward approach achieves high reconstruction speed, it still suffers from the substantial storage cost of 3D Gaussians. Existing 3DGS compression methods relying on scene-wise optimization are not applicable due to architectural incompatibilities.
%
To overcome this limitation, we propose TinySplat, a complete feedforward approach for generating compact 3D scene representations. Built upon standard feedforward 3DGS methods, TinySplat integrates a training-free compression framework that systematically eliminates key sources of redundancy. Specifically, we introduce View-Projection Transformation (VPT) to reduce geometric redundancy by projecting geometric parameters into a more compact space. We further present Visibility-Aware Basis Reduction (VABR), which mitigates perceptual redundancy by aligning feature energy along dominant viewing directions via basis transformation. Lastly, spatial redundancy is addressed through an off-the-shelf video codec.
% 为了解决上述问题，我们提出了一种新颖的压缩框架来消除Gaussian数据中不同形式的冗余。首先，我们提出VPT对几何参数进行坐标变换以消除结构冗余。然后，我们提出VABR通过计算一组球面基函数将编码能量集中在观察方向附近，消除感知冗余。最后，我们用一个通用视频codec来提取空间冗余。
% To fill this gap, we propose a training-free compression framework充分利用数据分布特性of Gaussians from the feedforward method. 我们提出了两种低复杂度变换方法来降低通道间相关性。具体来说，我们提出view-projection transform将高斯几何参数可逆地变换到相机坐标空间，where我们观察到数据的分布更加有规律。Moreover, 对于高维的球谐系数，我们提出一种球基函数变换。通过计算一组能量集中在观察方向附近的球面基代替SH基，从而获得更紧凑的色彩特征。我们进一步使用视频codec作为底层编码引擎来提取feature map中的空间冗余。
Comprehensive experimental results on multiple benchmark datasets demonstrate that TinySplat achieves over 100$\times$ compression for 3D Gaussian data generated by feedforward methods. Compared to the state-of-the-art compression approach, we achieve comparable quality with only 6\% of the storage size. Meanwhile, our compression framework requires only 25\% of the encoding time and 1\% of the decoding time.

% 近期，feedforward 3DGS的出现提供了一种三维场景重建的新范式。其通过在大规模数据集上训练神经网络来前馈生成3DGS模型from稀疏视点。由于这些特性，现有研究难以应用于前馈生成的3D高斯数据。本文中我们提出了一种Training-free compression Framework based on an existing video codec，专门针对上述场景进行优化。在我们的框架中，我们提出了view-projection transform和spherical basis tranform用去去除高维3DGS的inter-channel冗余。

% 现有的针对3D高斯数据的压缩研究主要集中在优化逐场景的3DGS训练策略。这些方法大多需要密集视点输入并且对每个三维场景进行独立优化。近期，feedforward 3DGS的出现提供了一条新的技术路线。其通过在大规模数据集上训练神经网络来前馈生成3DGS模型from稀疏视点。由于这些特性，现有研究难以应用于前馈生成的3D高斯数据。本文中我们提出了一种Training-free compression Framework based on an existing video codec，专门针对上述场景进行优化。在我们的框架中，我们提出了view-projection transform和spherical basis tranform用去去除高维3DGS的inter-channel冗余。
% 在多个common测试数据集上的全面的实验结果表明，我们的方法可以实现对此类数据超过100x的压缩。相比于SoTA方法，我们仅用不到10%的存储空间实现了同等的质量。同时我们的编码框架具有更低的计算复杂度。
\end{abstract}

\begin{IEEEkeywords}
3D Gaussian Splatting, Data Compression, Novel View Synthesis, Feedforward 3DGS.
\end{IEEEkeywords}

\section{Introduction}
\IEEEPARstart{T}{he}
3D Gaussian Splatting (3DGS) has emerged as a powerful technique for reconstructing 3D scenes from multi-view images. 
By explicitly representing a scene as a collection of anisotropic Gaussian primitives, 3DGS enables photorealistic rendering and supports real-time free-viewpoint navigation. 
The vanilla 3DGS formulation relies on stochastic gradient descent (SGD) and adaptive density control to jointly optimize the geometric and appearance attributes of the Gaussian primitives. 
While this approach achieves impressive visual fidelity, its reliance on per-scene optimization and dense input views significantly limits its scalability and practicality.

Recent feedforward 3DGS utilizes Neural Network~(NN) to directly infer the parameters of 3D Gaussian primitives from input images, thereby eliminating the need for iterative optimization. 
They offer substantial improvements in reconstruction speed and are capable of handling sparse view inputs, making them attractive for time-sensitive applications such as AR/VR and mobile 3D capture. 
However, the efficiency comes at the cost of significantly increased data volume. 
3DGS requires numerous Gaussian primitives to model a specific scene, resulting in substantial memory and storage demands.
This expansion poses major challenges for data transmission and interactive rendering, highlighting the need for effective compression techniques specialized for the unique characteristics of 3DGS data.

Prior research has explored optimization-based strategies for compressing 3DGS models, such as pruning~\cite{papantonakis2024reducing, lee2024compactrad, shi2024lapisgs, wang2024end2end, fan2024lightgaussian}, clustering~\cite{lu2024scaffold}, motion estimation~\cite{gao2024hicom, tang2025compressing}, and probabilistic modeling~\cite{chen2024hac, wang2024contextgs}. Notably, HAC~\cite{chen2024hac} employs multi-resolution hash grid priors for compact encoding, while ContextGS~\cite{wang2024contextgs} enhances compression efficiency through hierarchical context modeling. Nevertheless, due to their reliance on iterative optimization, these methods inherit the same limitations as vanilla 3DGS and are thus unsuitable for feedforward methods.

To overcome these limitations, FCGS~\cite{chen2024fcgc} introduced an optimization-free compression framework applicable to arbitrary 3D Gaussian models. By integrating hash grid-based hyperpriors and context priors through a multi-path entropy module, FCGS achieves notable compression ratios across diverse datasets. However, its effectiveness for feedforward-generated Gaussian models remains limited. First, FCGS does not fully exploit the inherent spatial redundancy derived from feedforward geometric generation. Additionally, the uniform treatment of feature channels hinders effective utilization of perceptual redundancy present in appearance features. Furthermore, reliance on multiple NN-based prior models introduces computational overhead unsuitable for real-time or edge-based scenarios.

In this paper, we propose \textbf{TinySplat}, a fully training-free approach for generating compact 3D Gaussian scene representations directly from multi-view images. TinySplat effectively addresses existing limitations by coupling a feedforward 3DGS generation stage with a rendering-aware compression stage. Specifically, in the Gaussian generation stage, we employ existing feedforward methods, such as DepthSplat~\cite{xu2024depthsplat}, to produce Gaussian feature maps, where each element defines parameters of individual Gaussian primitives. Subsequently, we introduce a novel compression framework designed to leverage spatial and perceptual redundancies within these feature maps for efficient storage and transmission.

In the compression framework, we propose \textit{View-Projection Transform} (\textbf{VPT}) to address structural redundancy from geometric generation. By exploiting the pixel-aligned characteristics of feedforward 3DGS, VPT reveals a more regular spatial layout, significantly enhancing local correlation and enabling more effective compression. Additionally, we propose the \textit{Visibility-Aware Basis Reduction} (\textbf{VABR}) method to handle perceptual redundancy in color attributes. Leveraging the anisotropic properties of spherical harmonic (SH) functions, VABR selectively retains perceptually dominant components while suppressing negligible ones, thereby improving compression efficiency without compromising rendering fidelity. Lastly, given the resemblance between Gaussian feature maps and traditional 2D images, we further exploit spatial redundancies using an off-the-shelf video codec.

Our main contributions are summarized as follows:
\begin{itemize}
    \item We introduce TinySplat, a fully feedforward pipeline consisting of Gaussian generation and a novel compression stage, enabling compact 3D Gaussian scene representations without scene-dependent optimization.

    \item We develop the VPT module, applying a coordinate-space transformation that exploits characteristics specific to feedforward geometric inference, thereby reducing structural redundancy and enhancing reconstruction quality.

    \item We propose the VABR module, which utilizes the visibility of SH bases to selectively preserve the most perceptually informative radiance components under constrained viewing directions, enabling highly compact color representations.

    \item Extensive experiments demonstrate that TinySplat can achieve comparable quality to DepthSplat with only 1\% of storage. In comparison to the state-of-the-art~(SoTA) compression method, TinySplat achieves a 90\% reduction in storage and a 75\% reduction in encoding time, simultaneously delivering superior visual fidelity..
\end{itemize}

The remainder of this paper is organized as follows. Section~\ref{sec:related works} provides an overview of related works on 3DGS generation and compression. Section~\ref{sec:method} describes the proposed TinySplat in detail. Section~\ref{sec:experiment} presents experimental evaluations and demonstrates the advantages of TinySplat. Section~\ref{sec:limitations} outlines limitations and future research directions. Finally, Section~\ref{sec:conclusion} concludes the paper.

% 从多视点图片/视频中恢复三维场景从而实时渲染自由视点视频可以带给用户前所未有的沉浸体验，从而深刻地变革音视频产业。

% 3D高斯泼溅技术在三维重建领域显示出了极大的潜力，激发了大量的相关工作。3D高斯泼溅技术基于逆渲染技术路线实现场景的三维重建，其在前人技术~\cite{}的基础上创新性的使用高斯椭球显式建模三维场景。vinilla 3D高斯泼溅技术通过SGD和自适应椭球分裂拟合场景的几何和属性信息，这种技术路线实现了高质量建模和实时自由视点漫游，Nevertheless其也有场景重建速度慢和输入视点要求密集的缺点。

% This approach offers a promising solution for efficient and high-quality 3D scene reconstruction, particularly in scenarios with limited input viewpoints.
% 最近，feedforward 3D高斯泼溅由于其出色的重建速度和稀疏视点重建能力逐渐引发了人们的关注。与先前基于优化的方法不同，feedforward 3D高斯泼溅直接使用神经网络推理从输入图片得到高斯椭球的位置和属性。从而这种方法具有接近实时的重建速度，同时有较好的稀疏视点合成效果。
% feedforward 3DGS估计深度图

% dispite 3DGS技术在渲染质量和速度上取得的进步，使用3D高斯泼溅进行3D场景重建的复杂度普遍很高，用户在一般消费级显卡上难以完成。为了解决这一问题，一个直观的解决方案是将服务端生成的3DGS数据传输到用户端，使用户得以在普通的硬件设备上进行自由视点观看。这就需要生成更加紧凑的3DGS模型。对于基于优化的3DGS泼溅，人们目前已经进行了广泛的研究以获得更紧凑的模型，包括剪枝，聚类，运动估计等。作为最高效方法的代表，HAC引入哈希网格先验实现了对3DGS数据的高效压缩，ContextGS进一步提出使用层次上下文先验进一步提高压缩率。

% 然而对于Feedforward 3DGS泼溅，目前压缩方法的研究仍然欠缺。FCGS是一种通用的Feedforward压缩方法，可以压缩任意的3D Gaussian模型in a training-free manner。FCGS精心设计了MEM对基于哈希网格的超先验和上下文先验信息进行融合。从而实现了对不同类型3D高斯数据超过20倍的压缩效率。

% 我们通过Gaussian Inference Network将输入图像转化成同分辨率的高斯 feature map，where每个像素保存了一个高斯基元的位置，形状和颜色参数。我们的编码框架被用于压缩上述feature map。
% training-free method面临一个问题，那就是参数失真并不代表。这导致了率失真优化的不一致，从而损伤了压缩效率。
% 我们提出使用off-the-shelf video codec作为backbone来去除feature map的空间冗余。基于, we propose View-Projection Transform~(VPT), projecting the original 3D Gaussian ellipsoids to input camera space。

% 本文中，我们注意到前馈3DGS推理模型的输出与V-PCC中的二维点云投影图具有很高的相似性。这启发我们使用基于视频codec的编码框架对其进行压缩。Further more,基于对3D高斯椭球几何参数统计特性的分析，我们提出对original 3D高斯进行view-projection transform，将其变换到相机空间进行编码以降低不同几何参数间的相关性。而对于表示色彩的球谐系数，我们注意到有限的观察视角对不同球谐系数重要性有决定性的影响。因而提出Spherical basis transform对SH进行快速的降维。

% 然而，这种方法仍然存在如下的一些问题，1、作为一种通用3D高斯压缩框架，FCGS并没有针对前馈方法的数据分布特点进行优化 2、3D高斯作为一种高维数据，其不同通道对渲染质量具有不同的重要性，FCGS并不能很好地处理这一问题 3、由于引入大量先验模型，FCGS需要进行大量串行的神经网络推理操作，速度较慢

% feedforward 3DGS通过NN可以将输入图像转化成同分辨率的高斯image，where每个像素保存了一个高斯基元的位置，形状和颜色参数。
%Feedforward 3DGS transforms input images into "Gaussian images" of the same resolution, where each pixel stores the position, shape, and color parameters of a corresponding Gaussian primitive. 

%Compared to the input pixel, the data volume of the Gaussian primitives increases by over two orders of magnitude. However, information theory dictates that applying a deterministic transformation to a source does not increase its information content. This indicates that the Gaussian representation contains substantial redundancy, not only in the spatial domain but also in dimensionality and precision.
%In this work, we propose a novel coding framework that effectively exploits spatial, dimensional, and precision redundancies to achieve efficient compression of Gaussian data. 
% 相比于输入像素值，高斯基元的数据量增加了上百倍。然而，信息论告诉我们对信源做函数映射并不会增加信息量。这indicate高斯数据中存在大量的冗余信息，不止是空间冗余，还包括维度冗余，精度冗余等。本文中，我们提出了一种novel coding framework来利用这些冗余实现对高斯数据的压缩。

% 单纯对3DGS进行逐通道的压缩存在一个limation，那就是编码时优化的失真是参数的失真，而对高斯的性能评价通常是靠衡量渲染视点与GT间的失真。这导致了率失真优化的不一致，从而损伤了压缩效率。
% 对于空间冗余，我们发现高斯feature map的空间分布特性和图像存在相似性。后者过去几十年已经得到了充分的研究，研究者开发了集成了预测、变换、后处理滤波等多种优化技术的混合编码框架。因而，我们使用现成的视频编码器作为我们框架中的底层编码引擎。
% 对于精度冗余，直观的想法是进行量化以降低精度。然而不同高斯基元的不同特征维度对于量化误差的敏感度都是不同的。例如，对离渲染视角不同距离的高斯点进行相同尺度的位移和缩放对渲染结果的影响显然是不同的。
% 这是出于如下观察，对离渲染视角不同距离的高斯点进行相同尺度的位移和缩放对渲染结果的影响是不同的。考虑到渲染新视角图像是在相机空间中进行的，直接在相机空间中进行编码可以有利于率失真优化。
% 多数的feedforward 3D高斯泼溅方法都利用深度图反投影操作生成与输入视点像素对齐的高斯点。这一特性使得此类3D高斯数据天然地适合类似V-PCC的编码框架。这启发我们提出了一种基于二维图像编码的编码框架用于压缩前馈式方法生成的3D高斯。然而，3D高斯数据与点云数据的一个显著差异是其具有维度更高的属性和几何特征，其中大部分是用于表达各向异性球谐系数。主成分分析降维是处理这种高维数据的有效手段，其通过最小化均方误差将原始数据投影到几个主方向上从而去除通道间冗余。主成分分析默认各个通道的数据同等重要，这对于球谐系数来说是不合理的，比如，高阶球谐和低阶球谐在特定查询方向上具有不同的缩放系数，另外与观察方向垂直的高阶球谐基可能几乎不会被看到，从而对结果不产生影响。这启发我们提出一种带权重的主成分分析降维方法，利用相机方向导出重要性先验以提高压缩率。

% 总的来说，我们的贡献可以被总结如下：
% 我们提出了一种基于视频codec的无需训练的编码框架，用于高效压缩基于前馈方法生成的3D高斯泼溅模型。
% 提出了view-projection transform作为预处理方法，以降低通道间冗余，从而提升压缩效率。同时利用SH函数各向异性的特性提出Spherical basis transform来对SH系数进行降维
% 结果显示我们的方法可以实现对于此类3D高斯模型xxx倍的压缩，同时PSNR损失在0.1dB以内。对比FCGS，我们的方法渲染性能损失更小，且压缩后的码流大小仅为FCGS的10%。同时编码速度比FCGS快10倍。

\section{Related Works}
\label{sec:related works}
\subsection{Novel View Synthesis via Per-scene Optimization}
Novel-view synthesis is a pivotal task in computer vision and graphics, aiming to generate photorealistic views from captured multi-view images.
Recently, per-scene optimization methods have achieved remarkable progress by formulating the 3D representation as learnable parameters optimized using SGD.
% 近年来，基于可微渲染的逐场景优化方法在该领域取得了显著进展。此类方法将三维模型作为待优化的参数，使用多视角图片作为训练数据集，通过渲染像素值和真实像素值计算损失函数，从而应用梯度下降技术对渲染模型进行渐进式的优化，最终获得高质量的三维模型。

As a seminal work in this field, the Neural Radiance Field~(NeRF)~\cite{mildenhall2021nerf} proposed by Mildenhall \textit{et al.} uses a deep multi-layer perceptron to learn color and opacity mappings. Novel view rendering is then performed through the volumetric rendering pipeline.
Muller \textit{et al.} further propose INGP~\cite{muller2022instant}, which models the scene using multi-resolution hash grids. By querying the explicit 3D structure, INGP can significantly reduce the computational cost of both training and inference.
Researchers have also explored various explicit data structures for scene representation, including voxel grids~\cite{sun2022direct, fridovich2022plenoxels}, multiple tensor planes~\cite{fridovich2023k, chen2022tensorf}, and octrees~\cite{yu2021plenoctrees}.
Based on NeRFs, Kerbl \textit{et al.} propose 3DGS~\cite{3dgs_vanilla}, which replaces volumetric rendering with a hardware-friendly rasterization technique. 3DGS enables real-time rendering on consumer devices and has inspired numerous follow-up works targeting different application domains~\cite{fu2024colmap, qian20243dgs, meng2024mirror, yu2024mip}. 
%Mildenhall等人提出了NeRF~\cite{mildenhall2021nerf},使用深层多层感知机学习每个位置的颜色和透明度。随后通过体积渲染进行场景的训练和渲染。进一步提出Instant-NGP，使用multiresolution hash grid建模场景，显著降低了训练和推理的计算开销。研究者们还提出了包括voxel grid, multiple tensor planes, octree在内的多种显式数据结构用于场景建模。在神经辐射场的基础上，Kerbl等人进一步提出3DGS~\cite{3dgs_vanilla}，使用硬件友好的光栅化技术代替体积渲染。3DGS实现了低复杂度实时渲染，激发了大量的后续研究工作~\cite{}。

% Nevertheless, 这些方法需要独立优化每个场景，因而有较大的训练开销。这给实际应用带来挑战。
% Neural Radiance Fields (NeRF)~\cite{mildenhall2021nerf} leverage a substantial Multilayer Perceptron (MLP) backbone and dense sampling strategies, significantly enhancing rendering quality but concurrently imposing computational overhead, resulting in slower synthesis processes. In contrast, 3D Gaussian splatting (3DGS)~\cite{3dgs_vanilla} delivers superior quality and faster rendering speeds by employing explicit representation. However, the primary limitation of 3DGS lies in its considerable storage demands and requirement for dense multi-view inputs to achieve precise scene reconstruction. This limitation critically restricts its deployment in edge computing environments where storage and computational resources are constrained. 

\subsection{Sparse View 3D Reconstruction}
Acquiring dozens of input views is usually impractical in real-world applications. As a result, researchers have investigated methods for 3D reconstruction from sparse input views. Several approaches~\cite{niemeyer2022regnerf, truong2023sparf, zhang2024cor, li2024dngaussian, han2024binocular} improve the 3DGS optimization process by introducing specialized regularization, such as view consistency and depth normalization. 
Meanwhile, some methods construct 3D scene representations using principles from multi-view stereo~\cite{pixelsplat, mvsnerf, chen2023explicit, szymanowicz2024flash3d, xu2024murf, mvsplat}. They utilize techniques like epipolar geometry and cost volume aggregation. Some methods also leverage Vision Transformer architectures to effectively fuse features across different views~\cite{noposplat}.
Further research also employs structural priors derived from large-scale diffusion models to maintain consistency between multiple viewpoints, enhancing the quality of novel view synthesis~\cite{gao2024cat3d, wu2024reconfusion, zhang2024gs, liu20243dgs, liu2024one, chen2024mvsplat360}. 
% Although these advances successfully address the challenges of dense-view scenarios, explicit representations used in novel-view synthesis still require considerable storage resources.

\subsection{feedforward 3D Gaussian Splatting}
Vanilla 3DGS~\cite{3dgs_vanilla} requires per-scene optimization, making it computationally intensive and time-consuming to obtain explicit scene representations. Recently, numerous feedforward 3DGS methods have been proposed, targeting 3D objects~\cite{szymanowicz2024splatter, tang2024lgm, zhou2024gpsplus, zheng2024gps, xu2024grm} or entire scenes~\cite{pixelsplat, noposplat, chen2024lara, szymanowicz2024flash3d, wewer2024latentsplat, min2024epipolar, chen2024mvsplat360, xu2024depthsplat, chen2024pref3r}. The feedforward methods focus on fast reconstruction of 3D scenes from sparse input views.
Specifically, these methods typically comprise a depth estimation module and a Gaussian attribute synthesis module. 
% The pipeline first estimates depth maps from sparse input views and performs back-projection to generate 3D Gaussian centers corresponding to each pixel. Then, the feature extraction and multi-view fusion networks infer geometry and color attributes for the 3D Gaussian ellipsoids.
% 具体来说，这类方法通常包含深度估计模块和高斯属性合成模块。前馈三维重建技术首先根据稀疏的输入视点进行深度估计，然后通过深度进行反投影得到一系列对应每个像素的高斯中心位置。最后，输入视点的原始图像经过特征提取和融合获得3D高斯的属性映射。

For single object reconstruction, GPS-Gaussian~\cite{zheng2024gps} and GPS-Gaussian+~\cite{zhou2024gpsplus} focus on 3D Gaussian reconstruction of the human body. Leveraging strong structural priors of human anatomy, these methods can generate a high-quality human body model from only a few input views. Szymanowicz \textit{et al.} propose Splatter Image~\cite{szymanowicz2024splatter}, which maps each input pixel to a 3D Gaussian using a simple yet efficient network, achieving real-time performance at 38 FPS for forward reconstruction. GRM~\cite{xu2024grm} further introduces a transformer-based architecture that effectively fuses multi-view information to achieve better reconstruction quality.
% 在单一物体重建方面，GPS-Gaussian~\cite{zheng2024gps}和GPS-Gaussian+~\cite{zhou2024gpsplus}专门面向人体的3D高斯重建任务，通过学习人体的结构先验信息，使用少量输入视点即可快速重建出高质量的人体模型。Szymanowicz等人提出了Splatter Image~\cite{szymanowicz2024splatter}，通过一个结构简单但高效的神经网络，将输入图像中的每个像素映射为一个3D高斯，从而实现了38 FPS的快速前向重建。GRM~\cite{xu2024grm}进一步引入transformer结构，能够高效融合多视图信息，将输入像素转换为像素对齐的三维高斯数据，并通过反投影获得三维结构。此类任务均受益于强结构先验，可以实现高精度、高保真的重建效果。
% Zheng等人和Zhou等人分别提出GPS-Gaussian和GPS++

In contrast, full-scene reconstruction poses greater challenges than single-object reconstruction due to increased scene complexity and lack of reliable structural priors. To address these challenges, Charatan \textit{et al.} propose PixelSplat~\cite{pixelsplat}, which leverages deep learning to predict dense probability distributions in 3D space and samples Gaussian centers accordingly to enable fast scene reconstruction. Chen \textit{et al.} introduce MVSplat~\cite{mvsplat}, which generates multiple depth maps per view and computes cross-view confidence scores to achieve more accurate geometric reconstruction.
Zhang et al. propose the Gaussian Graph Network (GGN)~\cite{zhang2024gaussian} to improve reconstruction quality with abundant input views. By constructing a graph structure, GGN enables each Gaussian primitive to aggregate features from multiple viewpoints, which significantly reduces the number of Gaussian primitives and effectively suppresses artifacts. 
% Zhang等人提出了高斯图神经网络~（Gaussian Graph Network, GGN）~\cite{zhang2024gaussian}用于提高输入视点较多的场景下的重建质量。GGN通过构建图结构使每个高斯基元都能融合多个视点下的特征，从而显著减少高斯椭球的数量并且降低了伪影。
Xu \textit{et al.} further propose DepthSplat~\cite{xu2024depthsplat}, leveraging pre-trained monocular depth features to improve both depth prediction and reconstruction quality, which achieves SoTA performance.
% 而对于完整场景，由于场景复杂性高、难以构建稳健的结构先验、稀疏视角下信息受限等问题，此类重建任务往往比单一物体的重建更加困难。Charatan等人提出pixelsplat~\cite{pixelsplat},通过深度学习预测三维空间中的密集概率分布，从中采样高斯中心位置，从而实现了场景的快速三维重建。Chen等人提出了MVSplat~\cite{mvsplat}，对每个视点生成一系列深度图，并通过夸视角匹配计算其置信度得分，从而实现更精确的几何重建。Xu等人进一步提出DepthSplat~\cite{xu2024DepthSplat}，构建了一个鲁棒的多视图深度模型，该模型利用预训练的单目深度特征同时提高深度估计和场景重建质量。验证了三维场景重建任务与深度估计任务之间的协同增益。


\subsection{3D Gaussian Splatting compression}
The substantial computational cost of generating 3D Gaussian models highlights the need for efficient storage and transmission, making compression a key challenge in 3DGS research.
Researchers have proposed numerous training-based methods to compress 3D Gaussian structures and reduce their storage and transmission costs.
% 这使得用户难以实时的构建3D高斯模型并进行自由视点漫游。直观的解决方案是直接存储和传输3D高斯模型而不是多视点图像数据。这就需要对其进行压缩。近期已经有很多针对基于训练的3DGS压缩方法。
% 从图片生成3D Gaussian 模型的巨大计算开销

Some works create sparser representations by pruning less significant Gaussian primitives during training~\cite{papantonakis2024reducing, lee2024compactrad, shi2024lapisgs, wang2024end2end, fan2024lightgaussian}. 
Vector quantization is also widely employed in 3D Gaussian compression~\cite{papantonakis2024reducing, lee2024compactrad, wang2024end2end, fan2024lightgaussian, navaneet2024compgs}, where a learnable codebook is used to quantize high-dimensional feature representations, achieving efficient compression. 
% 矢量量化同样被广泛用于对3D高斯的压缩，通过构建可学习的码表对高维的高斯特征进行量化，从而实现高达20x的压缩率。
Some other methods map 3D Gaussian primitives onto feature planes~\cite{morgenstern2024compact_selforg,lee2025compression} to compress with conventional video codecs. The feature planes require end-to-end optimization to obtain more compressible representations.
Further research focuses on generating more compact Gaussian models by meticulously designing 3D representations. Lu \textit{et al.} propose Scaffold-GS~\cite{lu2024scaffold}, clustering Gaussian primitives into anchor points with neural features.
More recent approaches utilize structural priors such as hash grids and context models to exploit the redundancies among 3DGS primitives, achieving about 100$\times$ compression compared to the vanilla 3DGS~\cite{chen2024hac,wang2024contextgs,chen2025hac++}.

%To further leverage contextual information in the encoding process,~\cite{chen2024hac, wang2024contextgs, chen2025hac++} improved compression efficiency by explicitly modeling the relationships between structured grids and Gaussians.
However, the aforementioned methods require an optimization process for each specific scene. They also require dense multi-view images as input. Both characteristics limit their applicability. To address this, Chen \textit{et al.} propose FCGS~\cite{chen2024fcgc}, a general-purpose feedforward compression framework. By incorporating multiple priors into the Gaussian mixture modeling, FCGS achieves efficient compression for arbitrary 3D Gaussian data, even outperforming many scene-specific optimization methods. Nevertheless, the generality of FCGS limits its ability to exploit the unique characteristics of different types of 3D Gaussian data, leaving room for further improvement in compression efficiency.
% 他们同时要求密集的多视点图片作为输入，这使得这些方法并不适用于feedforward 3DGS. 而针对feedforward 3DGS的压缩算法研究仍非常缺乏。Chen .et al 提出了FCGS~\cite{chen2024fcgc}，这是一种通用的前馈压缩框架。FCGS通过融合多种先验信息进行概率建模实现了对任意3D高斯数据的高效压缩，甚至超过了很多逐场景优化的方法。Nevertheless，FCGS的通用性也使其无法利用不同类型3D高斯数据的特性来进一步提升压缩效率。因此，在本文中我们提出了一种针对feedforward 3DGS的training free压缩框架。

%Recently, FCGS~\cite{chen2024fcgc} proposed a feedforward network that enables end-to-end compression of optimized primitives, although it still focuses on optimized 3D Gaussian primitives. In contrast, feedforward generative approaches inherently yield more compact 3D Gaussian representations than optimized primitives by leveraging feature maps.
%Based on this observation, we propose a method that directly compresses these feedforward generated compact representations, significantly enhancing compression efficiency and performance.

\begin{figure*}[t]
\centering
\includegraphics[width=0.98\textwidth]{imgs/framework.pdf}
\caption{The overall framework of TinySplat. We generate Gaussian feature maps from existing Gaussian inference networks in a feedforward manner and compress them with a meticulously designed compression framework. The symbols $\boldsymbol{\alpha}^v, \boldsymbol{r}^v, \boldsymbol{s}^v, \boldsymbol{\mu}^v, \boldsymbol{SH}^v$ denote the opacity, rotation, scale, mean position, and SH maps associated with view $v$, respectively. These parameters jointly define the geometry and appearance of the Gaussian primitives. In our compression framework, we first apply the VPT and VABR to reduce cross-channel redundancy. Subsequently, all features are quantized into 14-bit integers, and each feature plane is independently encoded as a grayscale image using the HEVC codec. On the decoder side, we perform dequantization and apply the inverse transforms to reconstruct the original Gaussian feature maps. Finally, we merge Gaussian maps from different views to form the complete 3D scene representation and render novel views.}
% 
\label{fig:framework}
\end{figure*}

\begin{figure}[t]
\centering
\includegraphics[width=0.98\linewidth]{imgs/VPT.pdf}
\caption{Illustration of the proposed VPT. We project the feedforward-generated Gaussian primitives from world space into the corresponding input camera space. We perform compression in camera space, reduce inter-channel correlation of geometric parameters for better efficiency. During decoding, the geometry is transformed back into world space to reconstruct the 3D scene.
}
\label{fig:vp-trans}
\end{figure}

\section{Method}
\label{sec:method}
% 图x 方法框架图
% 我们通过混合的特征提取网络，直接推理得到3DGS在特征域上的紧凑表达
% Transformer有xxxx优点，可以实现xxx
% 因此，我们参考Vision Transformer~\cite{}和DPT~\cite{}提取多视点图片的特征，将输入图片转换为隐向量表示
% 是否需要修改方法，目前视角2的F_vit由视角1的F_vit、视角1+2的F_conv经过CNN预测编码，是否需要对视角2的F_conv也添加视角间预测？可以解释为现在的方式能更好地学习运动信息以便于对数据量较大的F_vit进行精准预测；似乎需要实验支撑。
% 或者使用更精细的分层编码？性能可能会更好。factorized prior->F1_conv->F1_vit->F2_conv->F2_vit
% 图x 残差的热力图？残差分布与输入视点的关系？用于说明视点间参考的效果，无参考时能量分布->有参考时能量分布，与轮廓以及视差的关系，F_vit和F_conv间的相关性。

\subsection{Preliminary}
% 广受关注的3DGS使用一个高斯椭球集合{\mu, s, r, \sigma, SH}建模三维场景，其中xxx表示xxx.
% 表述高斯分布如何被建模成旋转和缩放
The widely recognized 3D Gaussian splatting employs a collection of Gaussian primitives to model 3D scenes. Each primitive consists of geometric and SH parameters that define its shape and radiance. Specifically, the geometric properties of Gaussian primitives are characterized by the Gaussian probability density function,
% 每个高斯球的几何特征由如下的概率密度函数刻画
\begin{equation}
    \mathcal{G}(\boldsymbol{x}) = \frac{1}{(2\pi)^\frac{3}{2}|\boldsymbol{\Sigma}|^{\frac{1}{2}}}exp(-\frac{1}{2}(\boldsymbol{x}-\boldsymbol{\mu})^T\boldsymbol{\Sigma}^{-1}(\boldsymbol{x}-\boldsymbol{\mu})),
\end{equation}
where $\boldsymbol{\mu}$ and $\boldsymbol{x}$ denote the Cartesian coordinates of the Gaussian center and the sample point, respectively. $\boldsymbol{\Sigma}$ is the 3D covariance matrix and $|\boldsymbol{\Sigma}|$ denotes its determinant. In the vanilla 3DGS, the covariance matrix is further decomposed into rotation matrix $\boldsymbol{R}$ and scaling matrix $\boldsymbol{S}$,
\begin{equation}
    \boldsymbol{\Sigma} = \boldsymbol{RSS}^T\boldsymbol{R}^T,
\end{equation}
where the diagonal scaling matrix $\boldsymbol{S}$ is stored as a vector $\boldsymbol{s}\in \mathbb{R}^3$ and $\boldsymbol{R}$ is stored as quaternion $\boldsymbol{q}\in \mathbb{R}^4$ that represents rotation. 

In the rendering process, each Gaussian primitive is projected into the camera space and integrated to produce a 2D density function, which is combined with the opacity parameter $\sigma$ to compute pixel-wise opacity. Meanwhile, the color value is obtained by querying SH functions with the view direction.
% 为了渲染一个给定视点，每个高斯椭球需要变换到相机坐标下，然后沿着视线方向积分，结合不透明度参数sigma计算二维的不透明度。颜色则通过球谐系数和view direction进行查询。整个渲染过程是完全可微的，因此可以通过与GT图像的差异计算损失梯度，从而优化场景的表达。
% In summary, the 3DGS modeling scenes with a series of Gaussian primitives, which are characterized by a collection of color and geometry parameters ($\mu$, s, q, $\sigma$, SH). 

To reconstruct a specific scene, the vanilla 3DGS approach first generates an initial set of Gaussian primitives via structure-from-motion techniques. Subsequently, this initial model undergoes optimization to match input ground-truth images, including adaptive densification guided by gradients to effectively capture fine geometric and textural details.
% 为了重建一个特定场景，the vanilla 3DGS通过sfm获取初始模型，然后通过随机梯度下降SGD优化初始模型以拟合输入GT图像，在这个过程中其使用梯度信息进行自适应的densification

Despite delivering high-quality reconstructions, this optimization-based approach suffers from slow processing speeds. To mitigate this limitation, recent studies propose eliminating the computationally intensive gradient descent step by introducing deep NNs enhanced with cross-view attention mechanisms. These networks trained on extensive multi-view datasets can directly produce accurate 3D scene representations through feedforward inference. To facilitate NN processing, this approach typically generates Gaussian feature maps of the same resolution as the input images, where each location stores all geometric and color attributes of a single Gaussian. Formally, the Gaussian inference process can be expressed as,
% despite很好的重建质量，这种方法的重建速度非常慢。因此研究者提出绕开消耗计算资源的梯度下降，使用神经网络进行前馈推理直接得到场景的三维表示。为了便于神经网络处理，这种技术路线通常生成与输入同尺寸的高斯椭球图，which可以表示为，
\begin{equation}
\label{feedforward gaussian manner}
    \begin{aligned}
    \mathcal{F}: \{( \boldsymbol{I}^v, \boldsymbol{K}^v,& \boldsymbol{E}^v)\}_{v=1}^V \Rightarrow \\
    &\{\cup(\sigma_i^v, \boldsymbol{\mu}_i^v, \boldsymbol{q}_i^v, \boldsymbol{s}_i^v, \boldsymbol{SH}_i^v)\}_{i=1,...,H\times W}^{v=1,...,V},
\end{aligned}
\end{equation}
% \textcolor{red}{figure?}
where $\boldsymbol{I}$ denotes the input image, $\boldsymbol{K}$, $\boldsymbol{E}$ denote the intrinsic and extrinsic parameters corresponding to the input image, respectively. $V$ denotes the number of viewpoints.
% Note that we are describing the most general scenario here; some recent methods may not rely on the extrinsics~\cite{ye2024no}.
% The inference network requires training on a large-scale multi-view image dataset. Given that both the inference and the 3DGS are differentiable, the discrepancy between the rendered images and the GT images can be directly backpropagated to optimize the inference network.
% 推理网络需要在大规模的多视点图像数据集上进行训练。考虑到网络的推理和3DGS都是可微分的，渲染图像与GT图像的差异可以直接反向传播计算推理网络的梯度。
%$s$ and $R$ denote the scale and rotation factor of the spatial Gaussian distribution, respectively. The $\sigma$ and $SHs$ represent the opacity and spherical harmonic coefficients characterizing anisotropic colors.

% 对3DGS质量的评估指标主要是其渲染成二维图像。
\subsection{Proposed Pipeline}
The overall pipeline of the proposed approach is illustrated in Fig.~\ref{fig:framework}. Initially, we employ pre-trained feedforward 3D Gaussian inference networks~\cite{mvsplat, xu2024depthsplat} to generate a set of 3D Gaussian primitives representing the target scene. For each input view, the network produces Gaussian primitives corresponding directly to image pixels, with each primitive consisting of geometric and color features as described in Eq.~\ref{feedforward gaussian manner}. These features collectively form per-view feature maps.

Compared to the original images, the size of this 3D Gaussian representation increases by more than two orders of magnitude. However, according to information theory principles, applying deterministic transformations does not inherently amplify the information content of the source data. Therefore, the generated 3D Gaussian model contains substantial redundancy, while existing compression methods for 3DGS relying on optimization are incompatible with the feedforward inference pipeline. To address this, we introduce a training-free compression framework designed to reduce both spatial and perceptual redundancies, producing a significantly more compact representation of the 3D scene.

Our compression framework comprises two lightweight yet highly effective modules, VPT and VABR. These modules specifically address the challenges posed by the absence of end-to-end optimization. First, high-dimensional Gaussian primitives inherently contain extensive redundancies. Second, without supervision from explicit distortion metrics, quantifying how deviations in parameters affect rendering quality becomes challenging.

The VPT targets geometric compression based on the following key observations. First, rendering occurs in camera space, making distortion measurements in this domain more directly representative of rendering quality impacts. Second, Gaussian positions derived via back-projection exhibit stronger spatial correlations in camera space. VPT leverages these properties by performing a reversible transformation of geometric features from world space into camera space, facilitating more efficient compression.

The VABR module addresses perceptual redundancy in appearance parameters by leveraging the anisotropic properties of SH functions. By assigning visibility-aware importance weights to SH coefficients and analyzing SH distributions across the whole scene, VABR derives perceptually consistent basis functions. This adaptive approach selectively retains color features that significantly contribute to dominant viewing directions, enabling efficient yet visually faithful compression.

% 由于渲染是在相机空间中进行的，我们进行VPT将高斯点的几何信息由世界坐标空间可逆地变换到相机坐标空间以降低通道间冗余。对于色彩参数，考虑到SH function各向异性的特点，我们提出VABR，根据观察方向来计算每个球谐系数的权重，从而尽可能保留对渲染结果贡献大的球谐系数。
%For the SH coefficients, which constitute the majority of Gaussian features, we adopt VABR to reduce the feature dimension. Given the anisotropic nature of SH functions, we take the viewing direction into account to compute weights for each SH feature dimension. This method can preferentially retain the coefficients that contribute most to the rendering result.
% 考虑到SH function各向异性的特点，我们考虑观察方向来计算每个球谐系数的权重，从而尽可能保留对渲染结果贡献大的球谐系数。

% 相比于输入像素值，高斯基元的数据量增加了上百倍。然而，信息论告诉我们对信源做函数映射并不会增加信息量。这indicate高斯数据中存在大量的冗余信息。
% 为了获得高斯数据更紧凑的表达，我们提出一种压缩框架，which can 挖掘空间冗余、精度冗余和维度冗余以实现对高斯数据的高效压缩。
% Due to the lack of end-to-end optimization, our training-free approach faces two primary challenges that limit compression efficiency. First, the high-dimensional Gaussian primitives exhibit significant inter-channel redundancy. Second, the absence of accurate distortion estimation during compression hinders effective rate-distortion trade-offs.
% To address these limitations, we propose two lightweight yet effective transformations: View-Projected Transformation (VPT) and Statistical Bitrate Reduction (VABR). These transformations significantly improve compression efficiency while introducing only marginal computational overhead.

% feature maps的总维度远高于图片，暗示着不同特征通道间可能存在很高的相关性。为了去除通道间冗余，我们提出了View-Projection Transform~(VPT)和Spherical Basis Transform~(VABR)。可以在仅引入marginal 复杂度的情况下提升压缩效率。
% 我们通过深度和相对偏移值表示点三维坐标，同时我们对高斯几何形状参数进行了近似的变换。实验表明通过坐标空间变换，可以使其分布更加集中，从而实现更好的压缩效率。
% 通过实验我们发现相比于世界坐标空间，将高斯椭球变换到相机空间后可以显著提高压缩率。且对于压缩失真更加不敏感，从而可以实现更小的失真。
Subsequently, we pre-quantize the transformed features into a 14-bit integer format based on their standard deviations. These quantized feature maps are then encoded as grayscale images using an off-the-shelf HEVC codec to further reduce spatial redundancy. Additionally, metadata required for reconstruction, including camera parameters for the inverse VPT, basis vectors for inverse VABR, and other dequantization-related information, is losslessly encoded and transmitted to the decoder.
% 另外，用于反VPT和VABR的相机参数，球基向量和用于反量化的元数据被无损编码并传输到解码端。

%In the Gaussian feature maps, we discover strong structural similarity to conventional 2D images, as shown in Fig.~\ref{}. This motivates us to propose a training-free compression framework that leverages general-purpose video codecs to encode these feature maps into compact bitstreams. This approach enables efficient storage and transmission without the need for additional training or domain-specific compression models. On the decoder side, the bitstreams are decoded to reconstruct the original Gaussian feature maps, which are then used to recover the 3D Gaussian scene representation for novel view synthesis.
On the decoder side, we first decode the feature maps using an HEVC decoder, followed by dequantization to recover the floating-point features. The reconstructed geometric parameters are then transformed back to the original world coordinate system through inverse VPT. Concurrently, the color feature coefficients are restored to their original SH representation using the received basis vectors. Finally, by integrating these decoded attributes, we reconstruct the complete 3D Gaussian representation, enabling the rendering of novel views using a standard 3DGS renderer.

% 在解码端，我们首先用HEVC解码器解码每个特征平面，然后进行反量化。之后，我们将重建的几何参数变换回世界坐标空间，同时将色彩特征系数恢复为原本维度的球谐系数。
% 通过重建的

% 3D高斯数据具有至少三十个通道，容易想到的是其中一些通道是高度相关的。为了降低通道间相关性，我们对输入数据的特征空间进行变换。对于几何参数，我们提出View-projection transform，对于占据大部分维度的球谐颜色特征，我们使用一种spherical basis transform在保证渲染质量的同时对其进行降维。随后，我们使用现有视频编码框架（H.265）分别编码每个通道。由于视频编解码器需要整型输入，我们对所有特征维度进行预量化将浮点数据转化为14bit的整数。
% 在解码侧，我们对视频解码器输出的重建数据进行反量化。然后进行逆变换以将高斯基元恢复到原本的特征空间。
% 最后，高斯拼接成

\begin{figure}[t]
\centering
\includegraphics[width=0.9\linewidth]{imgs/dist.pdf}
\caption{Statistical distributions of positions and scaling factors before and after VPT. After applying VPT, most of the positional energy is concentrated in the depth channel~(top row), and the distribution of scaling factors becomes more compact.
}
\label{fig:distribution}
\end{figure}

\subsection{View-projection Transform}

\label{sec_proj_trans}
The Gaussian models generated by the inference network typically exhibit pixel-wise alignment with the input views, distinguishing them from conventional 3D Gaussian representations. This characteristic introduces a highly structured spatial arrangement into the generated models. Specifically, each Gaussian center is computed by back-projection, resulting in an ordered distribution of 3D positions. Moreover, the projected areas of these Gaussian ellipsoids remain relatively uniform across the input view, causing the scale parameters to correlate strongly with their distances to the camera. This reflects a systematic depth-dependent scaling pattern.

% 高斯推理网络生成的高斯点是对齐到输入视点图像的每个像素的，这一特性引入了一些不同于其它3D高斯模型的性质。首先，高斯中心的位置是通过利用深度图进行反投影操作得到的，这使得高斯中心的三维坐标分布具有很强的规律性。其次，在渲染时每个高斯椭球需要在给定视点方向上进行投影，由于像素对齐的性质，每个椭球在输入视点上投影的区域大小通常是相近的。从而距离视点越远的高斯椭球，其尺寸参数也相应的越大。换句话说，这类3D高斯模型的缩放系数与深度之间存在较高的相关性。

Based on the above analysis, we propose to apply the VPT to the geometry parameters of Gaussian primitives, conditioned upon input camera parameters.
% 基于上述分析，我们提出根据视点坐标对高斯基元的位置和缩放参数进行View-Projection transform。这一变换来源于计算机图形学中渲染常用的MVP变换，其中，model transform将物体的局部坐标转换为世界坐标，而view and projection transform将三维世界坐标变换到相机坐标，where z轴为相机的朝向，从相机穿过每个像素的射线具有相同的x和y坐标，z坐标则表示点到相机的距离。
A conceptual illustration of the VPT is shown in Fig.~\ref{fig:vp-trans}, where we transform the Gaussian geometric parameters into the corresponding camera space.
Specifically, we first apply the view transformation to the position of Gaussian centers,
% 本文中我们发现将高斯点变换到相应视点的相机坐标空间可以有效降低高斯点形状特征间的冗余，从而提升编码效率。view-projection transform的过程如图~\ref{fig:vp-trans}所示。具体来说，我们对于坐标参数和旋转参数分别进行如下变换，
\begin{equation}
    z_{i,j}^v(x_{i,j}^v,y_{i,j}^v,1)^T = \boldsymbol{K}^v\cdot(\boldsymbol{R}^v\cdot \boldsymbol{\mu}_{i,j}^v + \boldsymbol{T}^v),
\end{equation}
where $\boldsymbol{K}^v$ denotes the intrinsic matrix of the input view $v$, while $\boldsymbol{R}^v$ and $\boldsymbol{T}^v$ are the rotation matrix and translation vector derived from the extrinsic parameters.
The vector $\boldsymbol{\mu}_{i,j}^v$ represents the Gaussian center in world coordinates. $z_{i,j}^v$, $x_{i,j}^v$, and $y_{i,j}^v$ represent the transformed position. In particular, $z_{i,j}^v$ corresponds to the depth map, while $x_{i,j}^v$ and $y_{i,j}^v$ represent normalized 2D coordinates in the image plane.
Given the pixel-aligned nature of feedforward-generated Gaussians, coordinates $(x_{i,j}^{v}, y_{i,j}^{v})$ closely cluster around their corresponding pixel centers $(\frac{i}{H}, \frac{j}{W})$, with $H$ and $W$ being the height and width of the input images, respectively. Using this regularity, we encode only the offsets relative to pixel centers to enhance compression efficiency.
% 其中K是相机的内参矩阵，R和T可由相机外参导出。Both are metadata used while generating per-view Gaussians，并且需要显式传输到解码端。$\mu_w$表示世界坐标。z, x, y为实际编码的参数，具体来说，z坐标的物理意义相当于深度图，而xy坐标则是图像平面上的归一化二维坐标。回顾之前提到前馈方法生成的3D高斯是像素对齐的，上述xy坐标在分布上非常靠近像素的二维坐标，如图~\ref{}所示，从而我们只需要对偏移值进行编码。

% 在严格意义上，椭球在透视变换中并不保持椭球的形状。因此对于椭球的形状参数，即rotation quantinions and scaling vectors，我们采取一种粗略的近似：
After transforming center positions into camera space, we approximate transformations for the Gaussian shape parameters. Since ellipsoids do not maintain exact geometric forms under perspective projection, we perform the following approximations,
\begin{equation}
    \hat{\boldsymbol{q}}_{i,j}^{v}=Quat(\boldsymbol{R}^v)\cdot \boldsymbol{q}_{i,j}^v,
\end{equation}
\begin{equation}
    \hat{\boldsymbol{s}}_{i,j}^{v}=f\frac{\boldsymbol{s}_{i,j}^{v}}{z_{i,j}^{v}},
\end{equation}
where $\hat{\boldsymbol{q}}_{i,j}^{v}$ and $\boldsymbol{q}_{i,j}^v$ denotes the rotation quaternions before and after the transformation, respectively, and $\hat{\boldsymbol{s}}_{i,j}^{v}$ and $\boldsymbol{s}_{i,j}^v$ denote the corresponding scaling factors. $Quat(\boldsymbol{R}^v)$ denotes the quaternion form of $\boldsymbol{R}^v$, $f$ represents the focal length.
The transformation of the rotation and scale parameters can be interpreted as rotating the Gaussian ellipsoids into the camera coordinate system, followed by a perspective scaling toward the focal plane.
Since all of these transformations are invertible, the approximations involved do not compromise rendering performance. 
% where a和b分别表示变换前后的旋转四元数，c和d表示变换前后的缩放系数。
% 由于上述变换全部是可逆的，这些近似并不影响渲染性能。

To illustrate the effectiveness of VPT, we analyze the distribution of Gaussian geometric parameters before and after the transformation, as depicted in Fig.~\ref{fig:distribution}. Compared to the initial distribution, the transformed Gaussian geometric parameters exhibit enhanced regularity and spatial coherence.

% 其中，xxxx。如图~\ref{}所示，不同于世界坐标中无序的分布。

\begin{figure}[t]
    \centering

    \subfloat{%
        \includegraphics[width=0.95\linewidth]{imgs/SHs30.pdf}%
        \label{fig:sh_30}%
    }

    \vspace{-0.5cm}
    
    \subfloat{%
        \includegraphics[width=0.95\linewidth]{imgs/SHs75.pdf}%
        \label{fig:sh_75}%
    }

    \vspace{-0.2cm}
    
    \caption{The variation of SH basis function values with respect to viewing direction. At a specific direction, a larger absolute value indicates that the corresponding SH coefficient has a greater influence on the final rendering result. $Y_l^m$ in the figure denotes the SH basis function of degree $l$, where $m$ ranges from $-l$ to $l$, representing $2l+1$ distinct spatial orientations.}
    \label{fig:sh_fig}
    % 球谐基函数的取值随观察视角的变化。在特定视角下，该值的绝对值越大说明相应维度的球谐系数越重要。$Y_l^m$表示l阶球谐基函数，which包含2l+1个不同的空间取向，对应m从-l到l的取值。
\end{figure}

\subsection{Visibility-aware Basis Reduction}
\label{sec_da_PCA}

% \begin{figure}[t]
% \centering
% \includegraphics[width=0.9\linewidth]{imgs/SBR_visualize.png}
% \caption{Visualization of the first six basis functions (left to right) with decreasing importance. Brightness encodes the absolute SH response in each direction relative to the input view.}
% \label{fig:sbr_visualize}
% \end{figure}

3DGS employs SH feature coefficients to model the anisotropic radiance distribution of each Gaussian ellipsoid. Although SH functions can effectively model the radiance with strong physical interpretability, from a compression perspective, the representation is notably redundant. Representing RGB colors using SH functions of order $N_l$ requires a total of $3 \cdot (N_l + 1)^2$ coefficients, leading to substantial storage overhead. However, Many SH basis functions correspond to spherical regions rarely or never observed from the available viewpoints, thus minimally contributing to rendering quality.

The redundancy is particularly pronounced in sparse view reconstruction tasks, where the range of novel views is limited due to the lack of visual information from distant viewpoints. 
As illustrated in Fig.~\ref{fig:sh_fig}, the response of SH functions varies significantly across viewing angles. Therefore, treating all SH coefficients equally without considering their visibility may reduce the compression efficiency.
% 球谐基作为球面函数的一组标准正交基，具有良好的物理可解释性。然而，从压缩的角度来看SH是极为冗余的。固定的基函数无法适应不同场景的几何特性，从而使得有些参数对结果几乎没有贡献。
%在稀疏视点重建任务中，需要合成的新视角通常是受限的。这是因为距离输入视角太远的视角由于缺少足够的输入信息而实际上没有意义。换句话说对于场景中的高斯基元，观测角度\theta \phi只在特定的区间上取值。这导致一些球谐基函数对渲染结果几乎没有贡献。如图~\ref{}所示，我们展示了SH基函数的取值随观测角度变化的情况。可以看到在不同角度区间内，球谐基的重要性变化很大。

To address this issue, we introduce the VABR, which linearly adapts standard SH bases into a compact, visibility-aware basis set. We first compute the visibility-based weights to capture the directional importance of each SH component. Each weight factor $\lambda_l^m$ is defined as the average absolute value of the corresponding SH basis function $Y_l^m$ over a region of valid directions,
% 首先，我们对于不同球谐基系数计算权重系数。权重系数$\lambda_l^m$可以定义为球谐基函数在有效视角区域内的平均绝对值：
\begin{equation}
    \lambda_l^m=\mathbb{E}_{\boldsymbol{D}}|Y_l^m(\boldsymbol{\omega})|=\frac{1}{|\boldsymbol{D}|}\int_{\boldsymbol{\omega} \in \boldsymbol{D}}|Y_l^m (\boldsymbol{\omega})|d\boldsymbol{\omega},
\end{equation}
where $l$ and $m$ represent the order and orientation factors of the SH function. $\boldsymbol{D}$ denotes the area of valid directions. In practice, it is challenging to rigorously define the area $\boldsymbol{D}$, so we approximate it via Monte Carlo integration,
\begin{equation}
    \lambda_l^m=\frac{1}{N_s}\sum_{i=1}^{N_s}|Y_l^m(\boldsymbol{\omega}_i)|,
\end{equation}
where $N_{s}$ is the number of sample directions.
% 然而，实际上难以精确地定义哪些方向是有效的。因此我们根据输入视点的参数导出一组方向样本以计算蒙特卡洛积分：
To obtain representative viewing directions, we cast camera rays through a $3 \times 3$ grid uniformly covering the entire image plane for each input view and aggregate them.
% 为了获得一组有效的视角方向，我们从图像平面中按照3x3网格从左上到右下采样9个点，\textit{i.e.} 左上角、正上方、右上角、...，以相机穿过每个点的射线方向作为对有效观察方向的采样

Subsequently, these weights inform the construction of a compact set of new spherical basis functions, which also reflect the statistical distribution of SH coefficients within the scene.
In particular, to ensure balanced perceptual contributions, we scale the original SH bases by their visibility-aware weights and linearly combine them into new optimized bases, resulting in the following transformation,
\begin{equation}
    \hat{\boldsymbol{Y}}(\boldsymbol{\omega}) = \boldsymbol{W}^T \cdot\boldsymbol{\Delta}^{-1}\cdot \boldsymbol{Y}(\boldsymbol{\omega}),
\end{equation}
where $\boldsymbol{Y}(\boldsymbol{\omega})$ and $\hat{\boldsymbol{Y}}(\boldsymbol{\omega})$ denote the original SH basis functions and the transformed basis functions, respectively. The weight matrix $\boldsymbol{\Delta}$ denotes the diagonal matrix consisting of $(\lambda_0^0, ...,\lambda_l^m,...)$. $\boldsymbol{W}\in \mathbb{R}^{d\times k}$ represents the transform matrix to be solved, which contains a set of d-dimensional unit orthogonal vectors, with $k$ representing the dimensionality of the transformed function space.
Subsequently, the forward and inverse transformations of the color coefficients are as follows,
\begin{equation}
    \boldsymbol{Z}=\boldsymbol{W}^{T}\cdot \boldsymbol{\Delta}\cdot  \boldsymbol{X},\ \ \hat{\boldsymbol{X}}=\boldsymbol{\Delta}^{-1}\cdot \boldsymbol{W}\cdot  \boldsymbol{Z},
\end{equation}
where $\boldsymbol{X},\hat{\boldsymbol{X}}\in \mathbb{R}^{d\times M}, \boldsymbol{Z}\in \mathbb{R}^{k\times M}$ represent the input, reconstructed, and transformed data matrices, respectively. Here, $M$ refers to the total number of Gaussian ellipsoids in the current scene.

To determine the transformation matrix $\boldsymbol{W}$ that can optimally preserve the perceptual information from higher-dimensional SH coefficients, we perform a statistical analysis on all SH coefficients in the current scene. Motivated by the formulation of principal component analysis, we compute the covariance matrix of $\boldsymbol{\Delta}\cdot \boldsymbol{X}$ and extract its principal components via eigen decomposition.
In particular, the eigenvectors corresponding to the $k$ largest eigenvalues are selected as the principal basis to form the transform matrix $\boldsymbol{W}$.

% 为了直观说明spherical basis transform的效果，我们对新的基函数进行了可视化， as shown in Fig.~\ref{}. 相比于球谐函数基在不同方向上同等重要，变换后的特征基更加侧重输入视点附近的观察方向。

% 获得逐通道的权重后，我们进行加权的PCA降维。具体来说，就是在执行原始的PCA降维前先将输入数据逐通道乘权重，在解码端将数据恢复到原始维度后再逐通道除权重。
% 我们在图~\ref{}中展示了逐通道权重系数的效果。降维过程可以看做是将SH函数进行线性组合

% Spherical basis transform对一个场景所有球谐特征进行分析和降维，从而达到数据压缩的目的。同时降维后的特征向量对量化误差倾向于具有相似的敏感度，有助于进行合理的码率分配。

% 我们认为球谐特征系数中可能存在大量通道间冗余。为了在编码时消除这种冗余，我们对每个场景进行PCA作为预处理。

% PCA默认所有输入维度同等重要，但是如表~\ref{}中展示的，在特定方向上每个球谐系数对渲染结果的贡献并不相同。另外，在实际应用场景中，需要合成的新视角是受限的。距离输入视角太远的视角由于缺少足够的输入信息而实际上没有意义。换句话说对于场景中的高斯基元，观测角度\theta \phi只在特定的区间上取值。这导致一些球谐基函数对渲染结果几乎没有贡献。



\begin{table*}[t]
  \renewcommand{\arraystretch}{1.3}
  \centering
  \caption{Quantitative comparison with 2 input views. TinySplat is compared with uncompressed models and FCGS on 3D Gaussians generated by different feedforward methods. ``Raw'' indicates the uncompressed baseline.}
    \begin{tabular}{cccc cc cc cc}
    \toprule
    \multirow{2}[1]{*}{\shortstack{Inference}} & \multirow{2}[1]{*}{\shortstack{Compression}} & \multicolumn{4}{c}{Re10k}     & \multicolumn{4}{c}{ACID} \\
\cmidrule{3-10}    & & PSNR(dB)$\uparrow$  & SSIM$\uparrow$  & LPIPS$\downarrow$ & Size(MB)$\downarrow$ & PSNR(dB)$\uparrow$  & SSIM$\uparrow$  & LPIPS$\downarrow$ & Size(MB)$\downarrow$ \\
    \midrule
    \multirow{2}[1]{*}{MVSplat} & Raw     & 26.36  & 0.8679  & 0.1291  & ~43   & 28.21  & 0.8419  & 0.1448  & ~43 \\
          & FCGS  & 25.90  & 0.8586  & 0.1533  & 2.650  & 27.61  & 0.8298  & 0.1800  & 2.596  \\
    (CVPR 2024) & Ours  & 26.32  & 0.8671  & 0.1349  & 0.216  & 28.13  & 0.8401  & 0.1528  & 0.229  \\
    \midrule
    \multirow{2}[1]{*}{DepthSplat} & Raw     & 27.50  & 0.8900  & 0.1125  & ~19   & 28.37  & 0.8481  & 0.1417  & ~19 \\
          & FCGS  & 27.35  & 0.8863  & 0.1201  & 2.674  & 28.27  & 0.8454  & 0.1466  & 2.666  \\
    (CVPR 2025) & Ours  & 27.43  & 0.8886  & 0.1192  & 0.166  & 28.29  & 0.8470  & 0.1501  & 0.181  \\
    \bottomrule
    \end{tabular}%
  \label{tab:main performance}%
\end{table*}%

\begin{table}[t]
  \centering
  \caption{$\alpha$ and $Q_{c}$ configurations for all data channels.}
    \begin{tabular}{ccccccc}
    \toprule
          & depth & offset xy & scale & rotation & color & opacity \\
    \midrule
    $N_{dim}$ & 1     & 2     & 3     & 4     & 6     & 1 \\
    \midrule
    $\alpha$ & 2048  & 256   & 256   & 256   & 1024  & 256 \\
    \midrule
    $Q_{c}$ & -4    & 12    & 0     & 9     & 3     & 0 \\
    \bottomrule
    \end{tabular}%
  \label{tab:qstep cfg}%
\end{table}%

\section{Experimental Results}
\label{sec:experiment}
\textbf{Datasets.} 
We conduct comprehensive experiments on several widely used benchmark datasets in the field. Among them, RealEstate10K~\cite{zhou2018stereo} and ACID~\cite{liu2021infinite} represent two large-scale multi-view 3D scene datasets that have been extensively adopted in related research. RealEstate10K dataset primarily comprises indoor scene multi-view images collected from online sources, containing 7,289 test scenes. The ACID dataset consists of numerous outdoor scenes captured by aerial drones, with 1,972 scenes for testing. In addition, we evaluate our method on a subset of the DL3DV dataset, containing 140 test scenes. Our test condition strictly adheres to the common configurations adopted in previous studies~\cite{pixelsplat, mvsplat, xu2024depthsplat}.
%我们在多个业界通用数据集上进行了实验。其中RealEstate10K和ACID是两个大规模的多视点三维场景数据集，被广泛应用于相关的研究。RealEstate10K主要包含从网络上获取的多个室内场景视频，其中包含67477个场景作为训练集，7289个场景作为测试集。ACID包含了无人机拍摄的大量户外场景，其中11075个场景被用作训练场景，1972个场景用作测试场景。另外，我们还使用了DL3DV数据集的一个子集，which包含140个测试场景。我们的测试条件完全follow先前的工作。

% 除去同时用于训练和测试的数据集，我们还在DTU数据集上测试了我们的预训练模型，以展示我们方法具有良好的泛化能力。DTU scan数据集包含一些物品不同视点的图片，which被广泛用于测试三维重建算法的性能

\textbf{Baseline.}
In the Gaussian generation stage, we adopt representative feedforward inference models, including MVSplat~\cite{mvsplat} and DepthSplat~\cite{xu2024depthsplat}, to generate 3D Gaussian feature maps. To demonstrate the generality and efficiency of the proposed pipeline, we compare our method with the optimization-free compression approach FCGS~\cite{chen2024fcgc} in terms of compression efficiency.
% We adopt the generic optimization-free 3D Gaussian compression approach FCGS~\cite{chen2024fcgc} as a baseline for the rate-distortion comparison.
% 为说明拟议压缩框架的通用性和高效性，我们将我们的方法基于若干代表性的Feedforward 3D Gaussian 推理方法进行了实现，包括pixelsplat~\cite{}, MVSplat~\cite{}, NoPoSplet~\cite{}. 为了进一步说明我们的框架具有很高的压缩效率，我们还引入了通用3D Gaussian压缩方法FCGS~\cite{}作为压缩效率的基线。我们将FCGS应用于不同的方法，然后提供在率失真尺度上的对比结果。

\textbf{Implementation details.}
In our experiments, the reserved color feature dimension $k$ for VABR is fixed at 6.
For 14-bit quantization, the quantization step for each channel is configured as $\sigma/\alpha$, where $\sigma$ represents the standard deviation of the current channel and $\alpha$ is an empirically chosen scaling factor, as detailed in Table~\ref{tab:qstep cfg}. 
We initially set $\alpha$ to 256 for each channel, and selectively increase it for a few error-sensitive channels to improve reconstruction quality.
% 我们将首先将quantization steps设置为256，仅对少数误差敏感的通道降低量化步长以提高质量。
Due to the varied distributions across different color basis components, we adopt a shared quantization step for all color channels to ensure consistent quantization. Here, $\sigma$ represents the standard deviation of the most dominant component.
% 特别的，由于色彩特征中对应各特征基的数据分布差异显著。我们对所有色彩特征使用相同的量化步长，where \sigma是最重要特征通道的标准差。
After quantization, an offset is applied so that the minimum quantized value becomes $0$.
Extreme outliers in each channel are truncated to ensure that the maximum quantized values do not exceed the valid 14-bit range.

In the video coding module, each channel is treated as an independent grayscale image and encoded using the HEVC reference software HTM15.0-RExt8.1, with all channels encoded in parallel via separate processes. To achieve different compression rates, we adjust a global Quantization Parameter~(QP), denoted as $Q_g$. Then, a set of per-channel QP offsets $Q_{c}$ is experimentally configured, as summarized in Table~\ref{tab:qstep cfg}. Subsequently, the QP values for the HEVC codec are configured as $Q_c+Q_g$.
All experiments are conducted on a workstation equipped with an Intel Core i7-13700K CPU and a single NVIDIA RTX 4090 graphics card, running in a WSL2.0 environment with Ubuntu 20.04.
% 实验中，我们将最终的球谐特征维度k设置为6。在14-bit量化部分，对于每个场景，我们将通道的量化步长设置为$\sigma/\alpha$，其中\sigma代表对应通道的标准差。我们实验性的将\alpha配置为如表~\ref{}所示。对于每个数据通道，我们对极端的离群值进行截断以保证量化后数据的动态范围不超过$2^14$。量化后的数据经过偏移处理以使最小的量化系数值为0。我们将每个通道视作独立的灰度图像使用HTM15.0-RExt8.1参考软件进行编码，每个通道在独立的进程中并行进行编码。我们使用一台搭载Intel core i7-13700K CPU和Nvidia RTX 4090的PC在WSL2.0+Ubuntu20.04环境下进行测试。
% 为了获得不同的压缩率，我们调整一个全局量化参数（QP），并实验性地为每个通道设置offset，如表所示


The overall syntax elements include Gaussian geometric parameters (position, scale, and rotation), color feature coefficients, and opacity.
In addition, a small amount of metadata is explicitly signaled to the decoder, including camera parameters, the basis vectors used for dimensionality reduction, as well as per-channel quantization steps and offset values.
% 码流中的语法元素包括高斯几何参数（位置，缩放，旋转），色彩特征系数，透明度。另外，还有一些占据很少空间的元数据需要显式signal解码器，包括相机参数、用于对降维的特征基向量矩阵、以及逐通道的量化参数和偏移值。

%我们将提出的压缩框架在MVSplat和depthsplat上进行了实现。将压缩模型的效果与原始模型以及FCGS进行对比。

\begin{figure*}[t]
\centering
\includegraphics[width=0.98\textwidth]{imgs/RD_all.pdf}
\caption{Rate-distortion performance on the Re10K dataset, where distortion is measured in terms of novel view synthesis quality. "Raw" denotes the uncompressed 3D Gaussian model from DepthSplat. Compared to the SoTA method, our method achieves comparable quality (PSNR, SSIM, and LPIPS) with only 6\% of the storage size.}
% Re10K数据集上的RD性能曲线。对比FCGS，我们的方法在同等质量（PSNR,SSIM,LPIPS）下码率仅为7\%。FCGS在PSNR尺度上下降更为明显
\label{fig:rd_vs_FCGS}
\end{figure*}

\begin{table}[t]
  \centering
  \caption{Compression performance under different numbers of input views. ``Raw'' indicates the uncompressed baseline generated via Depthsplat.}
  \vspace{0.5em}
  \footnotesize
  \begin{tabular}{cccccc
                }
    \toprule
    Method & Views & PSNR (dB)$\uparrow$ & SSIM$\uparrow$ & LPIPS$\downarrow$ & Size (MB)$\downarrow$ \\
    \midrule
    Raw   & \multirow{3}[2]{*}{2} & 19.62  & 0.6265  & 0.2986  & ~33 \\
    FCGS  &       & 19.26  & 0.6183  & 0.3061  & 4.518  \\
    Ours  &       & 19.52  & 0.6196  & 0.3147  & 0.615  \\
    \midrule
    Raw   & \multirow{3}[2]{*}{4} & 23.16  & 0.7792  & 0.1745  & ~66 \\
    FCGS  &       & 22.23  & 0.7619  & 0.1897  & 9.054  \\
    Ours  &       & 22.99  & 0.7694  & 0.1923  & 1.275  \\
    \midrule
    Raw   & \multirow{3}[2]{*}{6} & 24.16  & 0.8191  & 0.1456  & ~100 \\
    FCGS  &       & 22.74  & 0.7943  & 0.1672  & 13.561  \\
    Ours  &       & 23.95  & 0.8085  & 0.1633  & 1.892  \\
    \bottomrule
  \end{tabular}
  \label{tab:compression_views}
\end{table}

\subsection{Objective Results}

\textbf{Overall Compression Performance.}
We first generate the 3D Gaussian models using 
the official implementations and pre-trained weights MVSplat\footnote{\url{https://github.com/donydchen/mvsplat}}~(SHA: 1f5e5486) and DepthSplat\footnote{\url{https://github.com/cvg/depthsplat}}~(SHA: 175b17a6). Then we compress the generated models using our compression framework to evaluate its effectiveness. 

We evaluate the rendering PSNR, SSIM, LPIPS, and compressed model size on two standard datasets with 2 input views and 3 target views, as summarized in Table~\ref{tab:main performance}. We employ a near-lossless compression configuration in our TinySplat, where $Q_g$ is set to 0. For the SoTA method FCGS, we faithfully reproduce its performance using the official implementation\footnote{\url{https://github.com/YihangChen-ee/FCGS}} to ensure a fair comparison. 
The results show that TinySplat achieves over 100$\times$ compression on the 3D Gaussian data produced by the baseline, with negligible loss in rendering quality. Compared to the SoTA training-free method, our approach still achieves a 15$\times$ higher compression ratio, while providing comparable rendering quality. 
The rendering quality suffers a noticeable drop when using the FCGS pipeline to compress Gaussian models from MVSplat. This performance degradation could be attributed to the distinctive distribution of 3D Gaussian data generated by MVSplat, as illustrated in the FCGS paper~\cite{chen2024fcgc}.
We highlight that we evaluate the performance on the complete datasets following the general configuration and report the average results to ensure a comprehensive and reliable comparison.
% 对于FCGS，我们根据官方实现诚实地复现了性能。表格中对于MVSplat输出的3DGS FCGS渲染性能损失较大。我们分析可能的原因是MVSplat输出的3DGS数据分布比较特殊。需要强调表中的性能完整的ACID数据集上测试的平均结果。
% 我们根据official code生成了三种baseline方法的码流，然后使用我们的编码框架进行压缩以测试其有效性。
% 我们在两个常用数据集上，测试了近无损压缩配置下的PSNR,SSIM, Lpips和编码输出的码流大小，the results are shown in Table~\ref{}。结果显示我们的方法在几乎不损失渲染质量的情况下，对于baseline方法输出的3D高斯数据可以实现超过100倍的压缩。即使对比soTA的training-free压缩方法FCGS，我们的方法依然在渲染质量更高的情况下实现十倍的压缩率提升。
% 对比通用压缩方法，FCGS，无压缩方法
% 表x RealEstate10k+ACID+dl3dv

% 除了两个输入视点的配置
Beyond the two-view testing configuration, we also evaluate the compression performance under varying numbers of input views on the DL3DV dataset, as shown in Table~\ref{tab:compression_views}. We follow the testing configuration in DepthSplat~\cite{xu2024depthsplat} and reproduce the reported performance. The results demonstrate that our TinySplat generalizes well, as no significant degradation in compression quality is observed with more input views.
% 我们在DL3DV数据集上测试了不同数量输入视点情况下的压缩性能，如表~\ref{}所示。我们复现了DepthSplat的测试配置。结果说明我们的方法具有很好的泛化性，对于多输入视点的高斯数据压缩失真并未明显增加。

\textbf{Rate-Distortion Performance.}
% 图x 提出方法的RD曲线（对比Fast Feedforward 3DGS Compression+不同方法的RD-curve）对比更多通用压缩方法，
To evaluate the impact of compression distortion on the rendering quality of 3D Gaussian models, we compare the rate-distortion performance of our proposed method against FCGS, as shown in Fig.~\ref{fig:rd_vs_FCGS}. To reduce the evaluation cost, we only conduct experiments on the first 200 test scenes from the Re10K and ACID datasets. The results show that our method achieves comparable or even better PSNR, SSIM, and LPIPS performance compared to FCGS with only 6\% of its storage space.
% 为了测试我们方法中压缩失真对3D高斯模型渲染质量的影响，我们展示了proposed method 与FCGS对比的率失真曲线，如图~\ref{}。为了减少测试量，对于Re10K和ACID两个大规模数据集我们仅选取前200个测试场景进行实验。结果显示我们方法仅使用不到2%的空间实现了与FCGS相同或更好的PSNR, SSIM以及LPIPS质量。

\begin{table}[t]
  \centering
  \caption{Runtime analysis. Our TinySplat is decomposed into key components. }
    \begin{tabular}{cccccc}
    \toprule
    \multirow{2}[4]{*}{Method} & \multicolumn{4}{c}{Encoding Time~(s)} & Decoding  \\
\cmidrule{2-5}          & VPT  & VABR & HEVC  & Total & Time (s) \\
    \midrule
    FCGS  & - & - & - & 4.074    & 4.095 \\
    TinySplat  & 0.003 & 0.004 & 0.971 & \textbf{1.016} & \textbf{0.042} \\
    \bottomrule
    \end{tabular}%
  \label{tab:time}%
\end{table}%


\textbf{Runing time Analysis.}
% 图x 时间分析（编码/解码/前向推理）
To assess the computational cost of our TinySplat, we measure the average encoding and decoding times, as reported in Table~\ref{tab:time}. TinySplat encodes a scene in approximately 1 second and supports real-time decoding. In contrast, the SoTA approach, which heavily relies on NN–based probabilistic models, requires about 4$\times$ more time for encoding and over 100$\times$ more for decoding.

The primary computational overhead in our pipeline arises from the HEVC video encoder. However, we note that the HEVC codec in our implementation is the reference software, which prioritizes coding efficiency and correctness over running speed. In practice, significantly faster encoding and decoding can be achieved through engineering optimizations~\cite{heng2014highly} or dedicated hardware implementations~\cite{7100895, 6317152}.
% 为了分析我们方法的计算代价，我们统计了我们方法的编解码时间，如表~\ref{}所示。结果显示我们的方法编码一个场景仅需1s，解码时间可以达到实时。作为对比FCGS由于使用了大量神经网络概率模型编码需要的编解码时间分别是我们方法的4倍和100倍。我们的方法的计算开销主要来自视频编解码模块，然而需要注意的是我们使用的是HEVC的参考软件，运行速度并非其主要的考量指标。对于HEVC codec通过工程优化和设计专用芯片通常可以实现更高的编解码速度~\cite{}。

\begin{table}[t]
  \centering
  \caption{Bit allocation under different compression ratios.}
    \begin{tabular}{ccccccc}
    \toprule
    $Q_g$    & Size~(KB) & position & scale & rotation & color & opacity \\
    \midrule
    0     & 164.63  & 21.4\% & 17.5\% & 13.0\% & 34.6\% & 13.5\% \\
    3     & 122.29  & 21.0\% & 18.0\% & 11.5\% & 35.7\% & 13.7\% \\
    6     & 89.91  & 20.5\% & 18.8\% & 10.2\% & 36.7\% & 13.8\% \\
    12    & 47.90  & 20.0\% & 20.1\% & 8.4\% & 39.0\% & 12.6\% \\
    \bottomrule
    \end{tabular}%
  \label{tab:bit_allocation}%
\end{table}%

\begin{figure*}[t]
\centering
\includegraphics[width=\linewidth]{imgs/subjective.pdf}
\caption{Subjective results on Re10k dataset. We compress the 3D Gaussian models from DepthSplat with both FCGS and the proposed TinySplat. FCGS exhibits severe watermark-like artifacts at 11$\times$ compression ratio.
In contrast, our method achieves up to 100$\times$ compression while maintaining visually indistinguishable rendering quality from the uncompressed reference.}
\label{fig:subjective}
\end{figure*}


\begin{figure}[t]
\centering
\includegraphics[width=0.9\linewidth]{imgs/comp_wise_ablation.pdf}
\caption{Component-wise ablation analysis. We separately disable the VPT module and the VABR module to test their impact on coding efficiency. The PSNR is measured between the rendered novel view and the ground truth.}
\label{fig:ablation}
\end{figure}


\begin{figure*}[t]
\centering
\includegraphics[width=0.98\linewidth]{imgs/subjective6v.pdf}
\caption{Subjective results under 6 input views on DL3DV dataset. Our TinySplat achieves superior objective and perceptual quality using only 15\% of the storage compared to the SoTA compression method. FCGS exhibits noticeable compression artifacts even at relatively high bitrates.
}
\label{fig:subjective-6v}
\end{figure*}


\textbf{Bit Allocation.}
To better understand the bit allocation characteristics of our framework and guide future optimizations, we encode 200 scenes under various QP settings and analyze the average storage distribution across different components, as summarized in Table~\ref{tab:bit_allocation}. The results show that the relative bit allocation remains largely consistent across compression ratios, with color parameters consistently occupying the largest share of the bitrate.
% 各部分比特数占比（平均）？
% 我们使用提出的框架在不同QP下编码了200个场景，并统计了不同组分的平均存储空间占比，如表~\ref{}所示。结果显示我们的方法中比特分配随压缩率上升整体变化不大。色彩参数占据了大部分的空间


\subsection{Subjective Results}
We adopt DepthSplat as the representative Gaussian inference method and compare subjective compression artifacts using FCGS and our proposed framework, as shown in Fig.~\ref{fig:subjective}.
In particular, we adopt the highest compression ratio setting offered by FCGS, which still results in a storage size several times larger than that of TinySplat.
The rendering results show that FCGS easily introduces prominent watermark-like artifacts, whereas TinySplat maintains visually negligible degradation, even at higher compression ratios.
% 我们使用DepthSplat作为feedforward Gaussian inference method的代表，分别使用FCGS和我们的方法进行压缩来对比压缩主观损伤。我们使用了FCGS的最大压缩率配置进行对比，where存储占用仍然数倍于我们的方法。可以看到FCGS在提高压缩率后出现了大量的水印状伪影。而我们的方法即使压缩率远高于FCGS，主观损伤相比无压缩模型仍然难以察觉。

We also evaluate the subjective performance under 6 input views, as illustrated in Fig.~\ref{fig:subjective-6v}. The results demonstrate that increasing the number of input views does not noticeably degrade the perceptual quality of our method, highlighting its strong generalization capability. In contrast, FCGS produces severe artifacts even at relatively low compression rates.
% 我们也测试了6输入视点情况下的主观性能，如图~\ref{}所示。结果显示更多输入视点基本不会损伤我们方法的主观质量，展示了我们方法的泛化能力。作为对比，FCGS即使在较小的压缩率条件下仍然产生了严重的artifact。

\subsection{Ablation Studies}
% 表格4 消融，w/o VP trans, w/o D-PCA, w/o both
\textbf{Component-wise ablation analysis.} We conduct an ablation study by separately disabling the proposed techniques to assess their impact on the compression performance, as shown in Fig.~\ref{fig:ablation}. We utilize 200 scenes from the Re10k dataset for testing. Results demonstrate that VPT can effectively reduce the loss of geometric information under the same bitrate, thereby improving rendering quality.
Meanwhile, VABR enhances compression efficiency by concentrating the energy of SH coefficients, enabling significantly higher compression ratios.

% 我们分别绘制了开关各技术状态下的RD曲线，如图所示。结果显示关闭Sepherical Basis transform会导致压缩率大幅降低，说明通过基函数变换可以极大地减少通道间冗余。实验结果表明，VPT通过在相机坐标下进行编码，在相同码率下可以降低高斯几何和位置信息的损失，从而提升渲染质量。SBR则可以有效集中SH系数的能量，从而实现更高的压缩率。

% 图AI~LD

\textbf{Ablation on VABR.} 
We conducted an experimental analysis to investigate the impact of the retained feature dimensions on compression efficiency. Starting from the most dominant component, we progressively added additional components to evaluate how encoding performance evolves.
In addition, we computed the storage cost associated with each individual dimension. As shown in Fig.~\ref{fig:SBR_ablation}, the results indicate that after applying the VABR, the first feature dimension concentrates the majority of the information, accounting for the largest portion of the total bitrate. The information content in subsequent dimensions decreases gradually. The overall rendering performance remains nearly unchanged as the feature dimensionality exceeds 6. These results validate the effectiveness of our VABR, which concentrates most of the meaningful information in the first few dimensions.
% 我们通过实验分析保留的特征维度对压缩效率的影响。我们从最重要的维度开始依次添加维度从而测试编码性能变化情况。我们同时计算了每个维度的编码比特，结果如图~\ref{}所示。结果显示进行Spherical Basis Transform后，第一个特征维度集中了绝大部分信息，占用最多的编码比特。随后每个维度的信息量逐渐减少，特征维度超过六维后性能几乎保持不变。结果证明了我们球面基变换方法的有效性。


%\textbf{Comparison with point cloud compression frameworks.} To better prove the efficiency of our geometry coding strategy, we compare our method with Geometry-based Point Cloud Compression~(G-PCC).

\section{Limitations}
\label{sec:limitations}
In this work, we do not modify existing 3D Gaussian inference networks; instead, we focus on developing a general compression framework applicable to the Gaussian data they generated. While our method effectively produces compact 3D Gaussian representations, there remains a notable gap between our achieved compression ratios and theoretical limits. Since 3D Gaussians are derived entirely from a few input images, information theory suggests that their entropy should not surpass that of the inputs. However, even after compression, the size of the 3D Gaussian models still exceeds that of the original images compressed with conventional codecs.
We attribute this discrepancy primarily to the inference process itself, where complex NN models introduce significant prior knowledge into the generated 3D Gaussian representations. Developing lightweight methods to explicitly extract and leverage these priors presents a promising direction for future research. In addition, redesigning or fine-tuning inference networks to generate inherently more compressible representations represents another important avenue for exploration.

% 虽然我们的方法对于前馈3D高斯可以实现高效的压缩，目前压缩率距离其理论上限仍有很大差距。前馈3D高斯数据完全由若干输入图像推理得到，根据信息论我们可以知道这种3D高斯数据的信息量应当不大于输入图像。然而，经过我们的方法压缩后的3D高斯模型大小仍然高于直接应用通用视频编码器压缩输入图像。我们认为一个重要的原因是复杂的3D高斯推理网络引入了大量模型先验信息。如何使用低复杂度的方法提取这种先验信息，以及如何调整推理网络以生成更加易于压缩的数据仍需要深入研究。而本文中我们并没有对现有3DGS推理网络进行修改，而只关注研究对于这类数据的通用编码范式。

% 另外，我们认为不同视点的3DGS图之间应当存在一定的信息冗余，本文中我们对此仅做了初步的探索，并未提出有效的解决方案。
% \textcolor{blue}{view-prior?}

\section{Conclusion}
\label{sec:conclusion}
In this paper, we presented TinySplat, a fully feedforward approach for generating compact 3D Gaussian scene representations directly from multi-view images without per-scene optimization. Bridging the gap between fast 3D Gaussian generation and practical deployment, TinySplat integrates a feedforward inference stage with a novel rendering-aware compression stage.
For the compression stage, we proposed a dedicated framework with VPT and VABR modules to reduce structural and perceptual redundancy, respectively. Furthermore, we employed a standard video codec to eliminate spatial redundancy. Extensive experiments demonstrated that TinySplat significantly reduces storage requirements while maintaining high rendering quality. 

These findings offer an intuitive understanding of the redundancy characteristics in feedforward-generated 3D Gaussian data, providing practical guidance for future exploration and system design.
We believe TinySplat paves the way for scalable and deployable 3D scene representation, and sets a solid foundation for future research in real-time neural graphics systems.


\begin{figure}[t]
\centering
\includegraphics[width=0.98\linewidth]{imgs/Ablation_SBT.pdf}
\caption{Ablation study of the VABR. We present the variation of rendering PSNR with feature dimensionality increasing, along with the storage cost of each individual feature channel.
The results show that our method more effectively concentrates energy into the leading dimensions, enabling a more compact and efficient feature representation.}
\label{fig:SBR_ablation}
\end{figure}

\bibliographystyle{IEEEtran}
\bibliography{TCSVT}


\vfill

\end{document}


